\documentclass[11pt]{article}
\usepackage[T1]{fontenc}
\usepackage{lmodern}
\usepackage[margin=1.08in]{geometry}
\usepackage{amsmath,amssymb,amsthm,mathtools}
\usepackage{microtype}
\usepackage[numbers,sort&compress]{natbib}
\usepackage[colorlinks=true,linkcolor=blue,citecolor=blue,urlcolor=blue]{hyperref}
\usepackage{enumitem}
\usepackage{booktabs,longtable,array}
\hypersetup{pdftitle={The irrationality exponent of positive rational logarithms is 2},
  pdfauthor={Ryan Matthew Casper},
  pdfsubject={Rational logarithms, irrationality exponents, and Lean formalization}}
\setlist{topsep=4pt,itemsep=3pt}
\setlength{\emergencystretch}{2em}
\newtheorem{theorem}{Theorem}[section]
\newtheorem{proposition}[theorem]{Proposition}
\newtheorem{lemma}[theorem]{Lemma}
\newtheorem{corollary}[theorem]{Corollary}
\theoremstyle{definition}
\newtheorem{definition}[theorem]{Definition}
\theoremstyle{remark}
\newtheorem{remark}[theorem]{Remark}
\title{The irrationality exponent of positive rational logarithms is \(2\)}
\author{Ryan Matthew Casper}
\date{October 7, 2026}
\begin{document}
\maketitle
\begin{abstract}
We prove that the natural logarithm of every positive rational number other
than \(1\) is irrational and has irrationality exponent \(2\).
The argument extends OpenAI's weighted logarithmic interpolation and
determinant method to distinct exponential coordinates, incorporating the
arithmetic denominator and analytic radius costs of a general rational base.
The resulting theorem is formalized using the released \(\pi\) library.
It includes the individual logarithms of \(2\) and \(3\).
\end{abstract}
We extended OpenAI's weighted logarithmic interpolation and determinant method to positive rational logarithms. We proved the required interpolation result at distinct exponential coordinates, incorporated the additional arithmetic and analytic costs, and formalized the resulting exponent-two theorem using the released \(\pi\) library.

Here ``positive rational logarithms'' means natural logarithms of positive rational numbers other than 1. The conclusions below concern each logarithm individually. They do not assert an irrationality exponent for the quotient log 2 / log 3.

The contribution statement describes our extension. The inherited geometric comparison and determinant framework belong to the released OpenAI library. The exact source boundary is recorded in §9. This document is an exposition of a machine-checked argument, with several substantial inherited lemmas stated and attributed explicitly. It is not a claim that an external human referee has already reconstructed the proof.

\section{The theorem and the contradiction to be excluded}

\begin{theorem} If \(a\in\mathbb Q\), \(a>0\), and \(a\ne1\), then \(\xi=\log a\) is irrational and has irrationality exponent 2. More explicitly, for every real \(\nu>2\), there is an integer \(Q\ge2\) such that
\[
\left|\xi-\frac pq\right|\ge q^{-\nu}
\qquad(p\in\mathbb Z,\ q\in\mathbb N,\ q\ge Q).
\tag{1}
\]
The threshold may depend on \(a\) and \(\nu\); it is uniform in \(p,q\). The theorem includes \(0<a<1\), where the logarithm is negative.
\end{theorem}

Our convention for the irrationality exponent is the supremum of the exponents admitting infinitely many reduced rational approximations with error less than \(q^{-\nu}\). The formal source connects this convention to (1) and verifies boundedness of the relevant set before using its supremum.

We prove (1) by excluding arbitrarily large denominators with error at most \(q^{-\nu}\), for any fixed \(\nu>2\). This slightly stronger exclusion is convenient for choosing finitely many very widely separated denominators.

No irrationality assumption enters that construction. Irrationality follows afterwards: if \(\xi=u/v\) were rational, the unreduced fractions \(ku/(kv)\) would give exact approximations with arbitrarily large denominators, contradicting (1). Dirichlet's theorem then gives exponent at least 2, while (1) excludes every exponent above 2.

The result is an exponent statement. It does not supply an explicit threshold, bounded continued-fraction coefficients, or a lower bound \(c/q^2\) with fixed positive \(c\).

The upstream method and library are cited in \citep{openai2026pi}.

\section{Weighted logarithmic interpolation at distinct exponential coordinates}

We first describe the geometric input, because it is the point at which the extension changes the original centers. The geometry and interpolation are over \(\mathbb C\); the matrices used later have rational entries.

Let \(m,K\ge1\), and let
\[
W=(w_0,\ldots,w_m),\qquad V=(v_0,\ldots,v_m)
\]
be positive rational weights. Let \(\sigma>0\) be rational. Define the conservative comparison constant
\[
\mathcal C(m,\sigma)=
2(m+2)^{m+2}\left(1+\frac{m+2}{\sigma}\right)^{m+2}.
\]
We impose the following conditions:
\[
K(1+3\sigma)^{m+1}\frac{\prod_iw_i}{\prod_iv_i}<1,
\qquad
(1+\sigma)w_i<v_i\quad(1\le i\le m).
\tag{2}
\]
For equally sized subsets \(A,B\subseteq\{0,\ldots,m\}\), if the largest index in their symmetric difference is positive and belongs to \(A\), require
\[
\mathcal C(m,\sigma)\prod_{j\in B}v_j
<\prod_{i\in A}w_i.
\tag{3}
\]

The centers are
\[
b_j=(y_j,c_{j1},\ldots,c_{jm}),\quad 0\le j<K,
\]
with distinct nonzero \(y_j\). There is no requirement that the \(c_{ji}\) be distinct. Near each center use formal coordinates
\[
t=Y/y_j-1,\qquad
u_i=X_i-c_{ji}-\log(1+t).
\tag{4}
\]
A jet packet of height \(H\) consists of the coefficients of \(t^su^\beta\) with \(v_0s+\sum_i v_i\beta_i<H\).

\textbf{Interpolation theorem used here.} There is a positive integer \(R\), clearing the weight denominators, such that for every sufficiently large integer \(n\), polynomials of \(W\)-weighted degree at most \(H=nR\) map surjectively onto the packets of height \(H\) at all \(K\) centers.

The cofinal heights \(nR\) suffice for the determinant argument. We do not need, and do not claim as a checked intermediate theorem, surjectivity for every sufficiently large real height.

\subsection{A curve estimate}

For an irreducible algebraic curve \(C\) in the affine space with \(Y\ne0\), work on its smooth complete normalization. At every place \(P\), set
\[
D_W(C)=\sum_P
\max\left(0,\frac{\operatorname{pole}_P Y}{w_0},
\frac{\operatorname{pole}_P X_1}{w_1},\ldots,
\frac{\operatorname{pole}_P X_m}{w_m}\right).
\]
The notation means the sum of the displayed local maxima; identically zero coordinates contribute no poles. For each branch over a center define
\[
h_P=\min\left(\frac{\operatorname{ord}_P t}{v_0},
\frac{\operatorname{ord}_P u_1}{v_1},\ldots,
\frac{\operatorname{ord}_P u_m}{v_m}\right),\qquad
S(C)=\sum_{P\text{ over centers}}h_P.
\]
Orders of identically zero series may be infinite. On a nonconstant curve the minimum remains finite.

The needed estimate is
\[
D_W(C)\ge(1+\sigma)S(C).
\tag{5}
\]

Suppose instead that (5) fails. The volume inequality in (2) implies, by counting monomials and jet conditions, that for every sufficiently large \(N\) there is a nonzero polynomial \(F_N\) of \(W\)-degree at most \(N\), whose \(V\)-weighted order at every center is at least \((1+3\sigma)N\). This count needs no independence of the vanishing conditions.

Use the commuting polynomial differential frame
\[
D_0=Y\partial_Y+\sum_{i=1}^m\partial_{X_i},
\qquad D_i=\partial_{X_i}\quad(i>0).
\]
In the chart (4), these become \((1+t)\partial_t,\partial_{u_i}\). Differentiation preserves the \(W\)-degree bound and loses at most its assigned \(V\)-cost in weighted order.

Every derivative word of cost at most \(\sigma N\) therefore has weighted order at least \((1+2\sigma)N\). Its restriction to each branch has order at least \((1+2\sigma)Nh_P\), whereas its total pole degree is at most \(ND_W(C)\). The assumed failure of (5) forces all these derivatives to vanish identically on \(C\).

Consider common zero loci of derivatives at successive cost cutoffs
\[
r\delta_N,\qquad \delta_N=\sigma N/(m+2).
\]
Choose nested irreducible components containing the affine part of \(C\). Their dimensions range from 1 to \(m\); consecutive stages eventually share a component \(Z\). On this component, the defining equations from the earlier stage have degree at most \(N\), and their further derivatives through cost \(\delta_N\) vanish.

We now use the inherited \textbf{persistent normal comparison lemma}. For a coordinate normal basis indexed by \(A\), and a differential-frame normal basis indexed by \(B\), persistence gives
\[
\prod_{i\in A}w_i\le
\mathcal C(m,\sigma)\prod_{j\in B}v_j.
\tag{6}
\]
The lemma compares an upper bound from weighted Bézout with a lower bound from rectangular transverse jets. Its large-\(N\) requirement is uniform in the component: the relevant rectangular cutoffs are at least 2 once \(N\) is sufficiently large in terms of the fixed weights, dimension and \(\sigma\).

This is a substantial reused lemma, not a new result about arbitrary analytic frames. The formal applications are mapped in §9 to the original persistent-weight and transverse-slice sources.

If some positive-index coordinate direction is not tangent to \(Z\), take the largest such index \(i\). The coordinate normal basis can include \(i\), with all larger directions tangent. Equations (3) and (6) force every frame normal basis to include \(D_i\). All other frame vectors lie in
\[
\ker(dX_i-dY/Y).
\]
Their normal rank is one less than the codimension, so the tangent space of \(Z\) lies in this hyperplane. Thus \(dX_i=dY/Y\) on \(Z\).

If \(Y\) were nonconstant, restrict to an algebraic curve on which it remains nonconstant and pass to a smooth complete model. An exact differential \(dX_i\) has zero residues, while \(dY/Y\) has residue \(\operatorname{ord}_P Y\) at each zero or pole. A nonconstant rational function has such a place, giving a contradiction. Consequently \(Y\) is constant on \(Z\).

In the remaining case all positive-index coordinate directions are tangent. Since \(Z\) is proper, its tangent space is their span; therefore \(dY=0\), and characteristic zero again makes \(Y\) constant. In either case \(Y\) is constant on \(C\).

Distinct \(y_j\) now matter: a constant-\(Y\) curve meets at most one center. Some \(X_i\) is nonconstant on \(C\). At that center \(t=0\) and \(u_i=X_i-c_{ji}\), so the zero/pole equality for this function gives
\[
v_iS(C)\le w_iD_W(C).
\]
The coordinate inequality in (2) contradicts the assumed failure of (5). This proves the curve estimate, counting every branch, including branches over singular points.

\subsection{From curves to bounded-degree interpolation}

Choose \(R\) so that \(R/w_i,R/v_i\) are positive integers. Compactify by the monomials of \(W\)-degree at most \(R\), retaining the constant coordinate as the affine chart. If \(L=\mathcal O(1)\), then on every normalized affine curve,
\[
\deg L=R D_W(C).
\]
The constant and pure-coordinate monomials attain the local maxima defining this degree.

Replace the logarithm in (4) by a polynomial truncation whose first omitted term has \(t\)-weight larger than every relevant \(v_i\). This leaves the weighted filtration and branch contacts unchanged. At center \(j\), form the finite-colength ideal
\[
I_j=(t^{R/v_0},(u'_1)^{R/v_1},\ldots,(u'_m)^{R/v_m}).
\]
These ideals give a coherent ideal \(I\), equal to the unit ideal away from the centers.

On the ordinary blowup, write \(I\mathcal O=\mathcal O(-E)\) and \(A=\pi^*L\). The source constructs the actual curve models and proves, for noncontracted affine curves,
\[
A\cdot C=RD_W(C),\qquad E\cdot C=RS(C).
\]
Affine curves with \(Y\equiv0\) meet none of the prescribed centers, so \(S(C)=0\) and their curve estimate is automatic. Together with (5) for curves meeting \(Y\ne0\), this makes \(A-(1+\sigma)E\) nef. Boundary curves avoid \(E\); on contracted curves \(A\) has degree zero and \(-E\) is positive by relative ampleness.

For a sufficiently large integer \(r>1\), \(rA-E\) is ample. The exact identity
\[
A-E=
\frac{r-1}{r(1+\sigma)-1}\,[A-(1+\sigma)E]
+\frac{\sigma}{r(1+\sigma)-1}\,[rA-E]
\tag{7}
\]
has positive coefficients, and supplies ampleness together with a uniform positive curve margin. This step does not infer ampleness from strict nefness alone.

For all sufficiently large \(n\), the ordinary blowup has
\[
\pi_*\mathcal O(-nE)=I^n,\qquad R^q\pi_*\mathcal O(-nE)=0\quad(q>0).
\]
These are eventual identities for ordinary powers of \(I\); they are not assertions about integral closures or about every small \(n\). Serre vanishing for \(A-E\), followed by the projection formula and Leray, yields
\[
H^1(X,I^nL^n)=0.
\]
Sections therefore surject onto the local quotients. The monomials generating \(I_j^n\) have weighted order at least \(nR\), so this surjection prescribes all packets below \(H=nR\). Completion does not change these finite-colength quotients.

Finally, eventual polynomial generation for the projective embedding expresses sections of \(L^n\) as degree-\(n\) homogeneous polynomials in its coordinates. On the affine chart their \(W\)-degree is at most \(nR\). This last step is essential: a Chinese remainder theorem without a degree bound would not give our interpolation statement.

\section{The rational matrix associated to good approximations}

Fix \(a=n/d>0\) in lowest terms and \(\nu>2\). Under the contradiction hypothesis choose finitely many approximations
\[
r_i=p_i/q_i,\qquad |\xi-r_i|\le q_i^{-\nu},\quad 1\le i\le m,
\]
with very large, successively separated denominators. Put
\[
w_i=\lceil\log q_i\rceil,\qquad v_i=w_i/\theta
\]
for a rational \(0<\theta<1\), and take the centers
\[
y_j=a^j,\qquad c_{ji}=jr_i,\quad 0\le j<K.
\]
Positive powers of \(a\ne1\) are distinct, whether \(a>1\) or \(0<a<1\). Since \(\xi\ne0\), sufficiently large approximations also have \(p_i\ne0\).

Choose \(F_0>2/\theta\) and set
\[
T_i=\lceil F_0w_i/v_0\rceil,\qquad
G_i(t)=\sum_{1\le k<T_i}\frac{(-1)^{k+1}t^k}{k}.
\]
The omitted tail has weight \(T_iv_0\ge F_0w_i>v_i\), so replacing the formal logarithm by \(G_i\) preserves the packets.

Columns of the interpolation matrix are monomials \(Y^hX^\alpha\) with \(w_0h+w\alpha\le H\). Rows are \((j,s,\beta)\) with
\[
v_0s+w\beta/\theta<H.
\]
The actual entry is
\[
[t^su^\beta]\,
\bigl(a^j(1+t)\bigr)^h
\prod_i(jr_i+G_i(t)+u_i)^{\alpha_i}.
\tag{8}
\]
In particular, its factor \(a^{jh}\) is retained.

Interpolation gives full row rank over \(\mathbb C\) at cofinal heights \(H=nR\). Since the entries are rational, the resulting nonzero square minor is rational as well. Select such a minor \(\Delta_H\) containing all rows. Let \(M_H\) be its row count, and write
\[
b_H=\frac{\sum_{\text{rows}}w\beta}{M_HH}.
\]
Then \(0\le b_H\le\theta\), and lattice counting gives
\[
M_H\sim
\frac{K\theta^mH^{m+1}}{(m+1)!\,v_0\prod_iw_i}.
\tag{9}
\]
Both determinant estimates below concern this same minor and are uniform in the selected columns.

\section{The arithmetic lower bound, including the rational base}

Use the actual conservative denominator choices
\[
L_i=\operatorname{lcm}(1,\ldots,T_i),\qquad
D_H=\prod_iL_i^{\lfloor H/w_i\rfloor}.
\]
Multiply column \((h,\alpha)\) by \(d^{Kh}\prod_iq_i^{\alpha_i}\), row \((j,s,\beta)\) by \(\prod_iq_i^{-\beta_i}\), and every entry by \(D_H\).

If \(\beta\le\alpha\), the resulting entry is
\[
D_Hn^{jh}d^{(K-j)h}\binom{\alpha}{\beta}
[t^s](1+t)^h
\prod_i(jp_i+q_iG_i(t))^{\alpha_i-\beta_i}.
\]
Otherwise it is zero. The lcm factors clear the coefficient denominators, so the scaled nonzero determinant is an integer and has absolute value at least 1.

Set
\[
\Lambda=4+\log4,\qquad w_*=\min_iw_i.
\]
The formal lcm estimate is \(\log L_i\le\Lambda T_i\). From \(T_i\le F_0w_i/v_0+1\), the column degree bounds, and
\[
w_i-1\le\log q_i\le w_i,
\]
we obtain
\[
\frac{\log|\Delta_H|}{M_HH}\ge-(1-b_H)-E_{\rm ar},
\tag{10}
\]
where
\[
E_{\rm ar}=
\frac{\Lambda F_0m}{v_0}
+\Lambda\sum_i\frac1{w_i}
+\frac{\theta}{w_*}
+\frac{K\log d}{w_0}.
\tag{11}
\]
The denominator \(d\) of the rational base produces the final term. It vanishes for integer bases such as 2 and 3. We use \(K\), not a sharpened \(K-1\), and the constant \(\Lambda\) above throughout the parameter selection.

\section{The analytic upper bound and its two alternatives}

Let \(\rho=100\max(1,|\xi|)\). For a selected monomial \(P\) and multi-index \(b\), define the entire function
\[
f_{b,P}(z)=[u^b]P(e^z,z+u_1,\ldots,z+u_m).
\]
For \(P=Y^hX^\alpha\), this is
\[
\binom{\alpha}{b}e^{hz}z^{|\alpha|-|b|}
\]
when \(b\le\alpha\), and zero otherwise.

At center \(j\), put \(z_j(t)=j\xi+\log(1+t)\). Then \(e^{z_j(t)}=a^j(1+t)\), while the \(X_i\) arguments in (8) are
\[
z_j(t)+u_i+e_{ji}+\tau_i(t),\qquad
e_{ji}=j(r_i-\xi),\quad
\tau_i=G_i-\log(1+t).
\]
Expanding the \(X_i\) powers gives a finite linear combination of row vectors
\[
\bigl([t^\ell]f_{b,P}(z_j(t))\bigr)_P
\]
with \(b\ge\beta\) and \(wb\le H\). The scalar coefficients do not depend on the chosen column \(P\).

The approximation hypothesis gives
\[
|e_{ji}|\le K e^\nu e^{-\nu w_i}.
\]
Each tail begins at \(T_i\). A contributing term using tail multiplicities \(e_i\) satisfies \(\sum_i e_iT_i\le s\); hence \(\sum_iw_ie_i<H/F_0\). On \(|t|=1/2\), tail and binomial estimates give the scalar bound
\[
\exp\{-\nu w(b-\beta)+H E_{\rm tr}\},
\]
where
\[
E_{\rm tr}=
\frac{\nu}{F_0}+\frac{\log2}{v_0}
+\frac{\log4+\log(2K)+\nu}{w_*}.
\tag{12}
\]

On \(|z|\le\rho K\), the entire functions are bounded by
\[
\exp\left\{H\left(\frac{\rho K}{w_0}
+\frac{\log(2\rho K)}{w_*}\right)\right\}.
\]
Moreover, \(|z_j(t)|<\rho K/2\) on the coefficient circle. Expand in Taylor powers of \(z/(\rho K)\). If two rows have the same pair consisting of \(b\) and Taylor order, their determinant contribution is zero. Thus, among \(n_b\) rows with the same \(b\), the sum of Taylor orders is at least \(\binom{n_b}{2}\).

Absolute convergence permits the determinant expansion and summation. For each resulting test matrix \(B\), factorial bounds and the geometric Taylor sums yield
\[
|\det B|\le\exp\{-c\sum_bn_b^2+M_HH(E_{\rm hol}+o(1))\},
\qquad c=\log2/4.
\]
This estimate also covers singular test matrices. The collision saving is in the unnormalized determinant exponent; normalization divides it by \(M_HH\). The holomorphic cost is
\[
E_{\rm hol}=
\frac{\rho K}{w_0}+\frac{\log2}{v_0}
+\frac{\log(2\rho K)}{w_*},\qquad
E_{\rm an}=E_{\rm tr}+E_{\rm hol}.
\tag{13}
\]

Fix \(A>\theta\) and \(0<\eta<1\). Let \(N_A(H)\) count the multi-indices with \(wb\le AH\), and set
\[
c_H=\frac{c\eta^2M_H}{H N_A(H)}.
\]
If at least \(\eta M_H\) rows use such low indices, Cauchy--Schwarz gives \(\sum_bn_b^2\ge\eta^2M_H^2/N_A(H)\), producing the collision saving \(-c_H\). Otherwise more than \((1-\eta)M_H\) rows have \(wb>AH\), producing the approximation saving \(-\nu(A(1-\eta)-b_H)\).

The number \(Q_H\) of finite expansion choices for a single row is polynomial in \(H\) for fixed parameters; a bound is
\[
(\lfloor H\rfloor+1)^{2m}(\lfloor H/v_0\rfloor+1).
\]
The total number of determinant choice families is at most \(Q_H^{M_H}\), so its logarithm divided by \(M_HH\) is at most \(\log Q_H/H\), which tends to zero. The normalized factorial errors also tend to zero. We obtain
\[
\frac{\log|\Delta_H|}{M_HH}
\le E_{\rm an}+o(1)+
\max\{-c_H,-\nu(A(1-\eta)-b_H)\}.
\tag{14}
\]
Here all parameters and chosen approximations are fixed before \(H\) tends to infinity. Lattice asymptotics in (9) give
\[
c_H\longrightarrow
c_\infty=\frac{c\eta^2K\theta^m}{(m+1)v_0A^m}.
\tag{15}
\]

\section{Parameter selection without circularity}

Fix \(a,\nu\) first. Thus \(\rho\) and \(\delta=\log d\) are fixed.

Choose positive rational \(\theta,A,B,C\) satisfying
\[
0<\theta<A<B<1,\quad
\nu(A-\theta)>1-\theta,\quad
C>1,\quad CB<1,\quad B<C\theta<1.
\tag{16}
\]
For existence, choose rational \(b\in(1/2,1-1/\nu)\), then a small positive rational \(s\), and put \(\theta=1-s,A=1-bs\). For small \(s\), the gap in (16) is positive and \(A^2<\theta\). Choose \(C\in(A/\theta,1/A)\) and then \(B\in(A,\min(1/C,C\theta))\).

Choose a small rational \(\eta>0\) so that
\[
g=\nu(A(1-\eta)-\theta)-(1-\theta)>0,
\qquad \epsilon=\min(g/2,1/2).
\]
Take \(F_0>2/\theta\) with \(\nu/F_0<\epsilon/3\).

Now increase \(m\), setting
\[
K=\lfloor C^m\rfloor,\quad w_0=B^{-m},
\quad v_0=2K\theta^mw_0.
\]
The ambient interpolation mass is exactly \(1/2\). Also
\[
\frac K{w_0}\le(CB)^m\to0,\quad
v_0\sim2(C\theta/B)^m\to\infty,\quad
\frac m{v_0}\to0,
\]
while (15) becomes
\[
c_\infty=\frac{c\eta^2(B/A)^m}{2(m+1)}\to\infty.
\]
Choose one finite \(m\ge1\) for which
\[
\frac{\Lambda F_0m+2\log2}{v_0}
+\frac{(\rho+\delta)K}{w_0}<\epsilon/3,
\qquad c_\infty>2.
\tag{17}
\]
Fix these parameters.

Since \(K\theta^m<1\), choose small rational \(\sigma>0\) satisfying
\[
(1+3\sigma)^{m+1}/2<1,\quad
(1+3\sigma)^mK\theta^m<1,\quad
(1+\sigma)\theta<1.
\]
The middle condition is retained in the formal parameter package; the distinct-\(Y\) interpolation proof uses the ambient mass inequality.

Only now choose the finitely many approximations. Arbitrarily large denominators let us make all \(w_i\) large and successively separated, satisfying (3), as well as
\[
\Lambda\sum_i\frac1{w_i}
+\frac{\theta+\log4+\log(2K)+\nu+\log(2\rho K)}{w_*}
<\epsilon/3.
\tag{18}
\]
There are finitely many product comparisons. Since \(v_i=w_i/\theta\) for positive indices, common factors cancel and each remaining comparison is met by increasing the next denominator sufficiently. The weight-separation requirement and the uniform rectangular cutoff depend on the fixed weights and \(\sigma\), rather than the center locations. The eventual interpolation threshold may depend on the fully selected centers; these are fixed before \(H\) tends to infinity.

Equations (12), (17), and (18) give
\[
E_{\rm ar}+E_{\rm an}<\epsilon<g,\qquad
c_\infty>1+E_{\rm ar}+E_{\rm an}.
\tag{19}
\]
This order absorbs both costs introduced by a general rational base: \(\log d\) in the arithmetic estimate and \(\rho\) in the analytic radius.

\section{Comparing the same nonzero determinant}

For sufficiently large cofinal \(H=nR\), interpolation supplies the minor, so (10) and (14) both apply.

If the maximum in (14) is the collision alternative, compatibility requires
\[
c_H\le1-b_H+E_{\rm ar}+E_{\rm an}+o(1)
\le1+E_{\rm ar}+E_{\rm an}+o(1).
\]
This contradicts the strict limiting margin in (19).

For the other alternative compatibility requires
\[
\nu(A(1-\eta)-b_H)-(1-b_H)
\le E_{\rm ar}+E_{\rm an}+o(1).
\]
The left side is at least \(g\): it decreases with \(b_H\), and \(b_H\le\theta\). This contradicts the other margin in (19).

Both alternatives fail. Hence arbitrarily large approximations of the stated quality cannot exist, proving (1). The irrationality and Dirichlet argument from §1 completes the exponent-two conclusion.

\section{What the verification establishes}

The extension consists of 55 Lean modules on top of 869 unchanged modules from the released \(\pi\) library. The source manifest contains all 924 files. Endpoint and intermediate axiom audits report only the standard foundational axioms: propositional extensionality, classical choice, and quotient soundness.

The proof export was accepted by Nanoda 0.4.19 and by a separate stock Lean 4.34.0 kernel, in addition to the original Lean 4.34.1 build. The second Lean check replayed the exported declarations into a fresh environment rather than trying to load newer binary objects into an older runtime. Negative controls checked rejection of a forbidden axiom and an ill-typed purported proof. Exact pins, commands, hashes, and limits are in \href{run:lean/VERIFICATION.md}{VERIFICATION.md}.

These checks validate the formal declarations relative to their definitions and foundational axioms. The human exposition still needs mathematical scrutiny, particularly at the inherited comparison lemma, the passage from curve inequalities to ordinary-power interpolation, and the determinant expansion. A separate statement check tests the connection between the formal declarations and the advertised approximation bound.

An AI reviewer with no proof history first formalized the advertised claim using only Mathlib definitions. We then proved its exact frozen type in \href{run:review/blind-statement/CertifiedBridge.lean}{CertifiedBridge.lean}. That separate Lean compilation succeeded and its axiom audit reported only the three foundational axioms above. The bridge does not import the specification file's intentional placeholder. This checks statement agreement; it is not a blind mathematical reproof.

This distinction addresses the proof-correctness and mathematical-understanding stages discussed in \href{https://arxiv.org/html/2608.16753v1}{Tao's essay}: kernel acceptance supports correctness of the formal statement, while readable arguments and independent mathematical scrutiny support understanding and wider acceptance.

\section{Source map and provenance}

The inherited OpenAI library is pinned to \href{https://github.com/openai/math/tree/adc7f1241b42e322a6451854ab7e4b4c146bf78a}{commit adc7f1241b42e322a6451854ab7e4b4c146bf78a}. Mathlib is pinned to d13f23b723b8a846827a245b89c10fc7d3f11612. The extension's source manifest and reproduction instructions are included in the \href{run:lean/README.md}{certified bundle}.

\small
\begin{longtable}{@{}>{\raggedright\arraybackslash}p{0.29\textwidth}>{\raggedright\arraybackslash}p{0.66\textwidth}@{}}
\toprule
\textbf{Mathematical step} & \textbf{Principal formal sources}\\
\midrule
\endhead
Weighted frames and generalized centers & \href{run:lean/Logarithm/GeneralizedCenters.lean}{GeneralizedCenters}, \href{run:lean/Logarithm/GeometryData.lean}{GeometryData}\\[5pt]
Auxiliary polynomial and derivative vanishing & \href{run:lean/Logarithm/AuxiliaryJet.lean}{AuxiliaryJet}, \href{run:lean/Logarithm/DerivativeVanishing.lean}{DerivativeVanishing}\\[5pt]
Inherited persistent normal comparison & \href{run:lean/OAI/NumberTheory/PiExponent/Approximation/PersistentWeightComparison.lean}{PersistentWeightComparison}, \href{run:lean/OAI/NumberTheory/PiExponent/Jets/TransverseSliceComparison.lean}{TransverseSliceComparison}\\[5pt]
Rigidity, constant fibers, and curve inequality & \href{run:lean/Logarithm/CurveRigidity.lean}{CurveRigidity}, \href{run:lean/Logarithm/ConstantFiber.lean}{ConstantFiber}, \href{run:lean/Logarithm/CurveInequalityIntrinsic.lean}{CurveInequalityIntrinsic}\\[5pt]
Actual degree identities and ampleness & \href{run:lean/Logarithm/CurveModelDegrees.lean}{CurveModelDegrees}, \href{run:lean/Logarithm/GeometryAmple.lean}{GeometryAmple}\\[5pt]
Ordinary powers and interpolation & \href{run:lean/Logarithm/CompactJetIdeal.lean}{CompactJetIdeal}, \href{run:lean/Logarithm/GeometricInterpolation.lean}{GeometricInterpolation}, \href{run:lean/Logarithm/Interpolation.lean}{Interpolation}\\[5pt]
Actual rational matrix and rank & \href{run:lean/Logarithm/ScaledMatrix.lean}{ScaledMatrix}, \href{run:lean/Logarithm/GlobalMatrixInterpolation.lean}{GlobalMatrixInterpolation}\\[5pt]
Arithmetic denominator bound & \href{run:lean/Logarithm/ScaledArithmetic.lean}{ScaledArithmetic}\\[5pt]
Analytic expansion and collision bound & \href{run:lean/Logarithm/LiteralAnalyticSummand.lean}{LiteralAnalyticSummand}, \href{run:lean/Logarithm/ScaledPeriodAnalytic.lean}{ScaledPeriodAnalytic}, \href{run:lean/Logarithm/AnalyticAggregate.lean}{AnalyticAggregate}\\[5pt]
Parameter existence and contradiction & \href{run:lean/Logarithm/LogarithmParameters.lean}{LogarithmParameters}, \href{run:lean/Logarithm/DeterminantContradiction.lean}{DeterminantContradiction}\\[5pt]
Irrationality exponent consequence & \href{run:lean/Logarithm/ExponentConsequence.lean}{ExponentConsequence}, \href{run:lean/Logarithm/Main.lean}{Main}\\[5pt]
\bottomrule
\end{longtable}
\normalsize

\textbf{AI use disclosure.} The author initiated and directed an investigation of extending OpenAI's irrationality-exponent method to logarithms of positive rational numbers. OpenAI coding agents contributed substantially to developing the candidate argument, producing the Lean formalization, and conducting automated adversarial reviews. The construction builds on OpenAI's published manuscript and released Lean library.

OpenAI agents also assisted in preparing this exposition and the blind specification check. These automated checks do not constitute external human peer review.

\renewcommand{\bibfont}{\raggedright}
\bibliographystyle{plainnat}
\bibliography{references}
\end{document}
